\documentclass[11pt]{article}

\usepackage[margin=1in]{geometry}
\usepackage{amsmath,amssymb,amsthm}
\usepackage{graphicx}
\usepackage[hidelinks]{hyperref}

\newtheorem{theorem}{Theorem}
\newtheorem{lemma}{Lemma}
\newtheorem{corollary}{Corollary}
\newtheorem{remark}{Remark}

\newcommand{\ellone}{\ell_1}
\newcommand{\ellinf}{\ell_\infty}
\newcommand{\HL}{\mathcal{H}L}
\newcommand{\G}{\mathcal{G}}
\newcommand{\N}{\mathbb{N}}

\title{A Full Answer to the \(\ell_1\)-Copy Question in arXiv:2605.21279}
\author{Automatic Math Discovery Run \texttt{fa\_banach\_001}}
\date{\today}

\begin{document}
\maketitle

\section*{Source}

Ver\'onica Dimant, Luis C. Garc\'ia-Lirola, Juan Guerrero-Viu, and
Anton\'in Proch\'azka, \emph{Composition operators for holomorphic Lipschitz
functions}, arXiv:2605.21279.

The question appears in Remark 4.21 on page 25 of the source PDF.

\begin{center}
\includegraphics[width=\textwidth]{figures/open_problem_crop.png}
\end{center}

\section*{Problem Statement}

For a holomorphic Lipschitz symbol \(\phi\), the paper writes
\[
  C_\phi^L=\widehat{\phi}^{\,*},
\]
where
\[
  \widehat{\phi}:\G_0(B_Y)\longrightarrow \G_0(B_X)
\]
is the corresponding holomorphic-Lipschitz linearization. Remark 4.21 asks:
if \(\widehat{\phi}\) fixes a copy of \(\ellone\), must it fix a complemented
copy of \(\ellone\)?

\section*{Result}

\begin{theorem}
Let \(T:Y\to X\) be a bounded linear operator between Banach spaces. If \(T\)
fixes a copy of \(\ellone\), then \(T\) fixes a complemented copy of
\(\ellone\).
\end{theorem}

\begin{proof}
Assume that \(T\) fixes a copy of \(\ellone\). Then there is an isomorphic
embedding
\[
  A:\ellone\longrightarrow Y
\]
such that \(TA:\ellone\to X\) is also an isomorphic embedding.

Consider the adjoint
\[
  Q=(TA)^*:X^*\longrightarrow (\ellone)^*=\ellinf.
\]
This map is onto. Indeed, a functional on \(\ellone\) defines a bounded
functional on the subspace \((TA)(\ellone)\subseteq X\) via
\((TA)^{-1}\), and Hahn--Banach extends it to all of \(X\).

Since \(\ellinf\) is injective, the quotient map \(Q\) has a bounded linear
right inverse. Thus there is a bounded operator
\[
  J:\ellinf\longrightarrow X^*
\]
such that \(QJ=I_{\ellinf}\).

Restrict \(J\) to the canonical copy of \(c_0\subseteq\ellinf\). For every
 \(u\in c_0\),
\[
  u=QJu=(TA)^*Ju=A^*T^*Ju.
\]
It follows that \(T^*J|_{c_0}:c_0\to Y^*\) is an isomorphic embedding: the
upper estimate is immediate from boundedness, while the lower estimate follows
from
\[
  \|u\| \leq \|A^*\|\,\|T^*Ju\|.
\]
Therefore \(T^*:X^*\to Y^*\) fixes a copy of \(c_0\).

Lemma 2.5 of the source paper states that, if \(T^*\) fixes a copy of \(c_0\),
then \(T\) fixes a complemented copy of \(\ellone\). Applying that lemma gives
the desired complemented copy.
\end{proof}

\begin{corollary}
Let \(X,Y\) be complex Banach spaces and let
\(\phi\in \HL_0(B_Y,X)\) satisfy \(\phi(B_Y)\subseteq B_X\). If the
holomorphic-Lipschitz linearization
\[
  \widehat{\phi}:\G_0(B_Y)\to \G_0(B_X)
\]
fixes a copy of \(\ellone\), then \(\widehat{\phi}\) fixes a complemented copy
of \(\ellone\).
\end{corollary}

\begin{proof}
Apply the theorem to \(T=\widehat{\phi}\).
\end{proof}

\section*{Verification Notes}

The proof uses three standard ingredients.

\begin{itemize}
\item If \(TA:\ellone\to X\) is an isomorphic embedding, then
  \((TA)^*:X^*\to\ellinf\) is onto by Hahn--Banach.
\item The space \(\ellinf\) is injective, so this quotient map splits.
\item Lemma 2.5 of arXiv:2605.21279 says that \(T^*\) fixing \(c_0\) implies
  \(T\) fixes a complemented copy of \(\ellone\).
\end{itemize}

No computational verification is relevant.

\section*{Limitations and Novelty Check}

The argument is a general operator-theoretic observation and does not use any
special geometry of holomorphic-Lipschitz free spaces. For that reason, the
result should be reviewed as a short deduction from classical facts together
with Lemma 2.5 of the source paper.

\section*{Verdict}

The question in Remark 4.21 has an affirmative answer as stated:
\(\widehat{\phi}\) fixing a copy of \(\ellone\) automatically implies that
\(\widehat{\phi}\) fixes a complemented copy of \(\ellone\).

\end{document}
